% Editable Beamer draft, 12 slides. Compile twice with pdflatex.
% Speaker notes (\note) are editable in this source; hidden in the PDF by default.
% To print notes, change "hide notes" below to "show notes".
\documentclass[aspectratio=169]{beamer}
\usepackage[T1]{fontenc}
\usepackage[utf8]{inputenc}
\usepackage{lmodern,amsmath,amssymb,booktabs,array}
\geometry{paperwidth=13.333in,paperheight=7.5in}
\definecolor{teal}{HTML}{087E83}
\definecolor{ink}{HTML}{193D42}
\definecolor{muted}{HTML}{58696C}
\setbeamercolor{normal text}{fg=ink,bg=white}
\setbeamercolor{frametitle}{fg=ink}
\setbeamercolor{structure}{fg=teal}
\setbeamercolor{alerted text}{fg=teal}
\setbeamerfont{normal text}{size=\fontsize{23}{31}\selectfont}
\setbeamerfont{frametitle}{size=\fontsize{34}{40}\selectfont,series=\bfseries}
\setbeamerfont{footline}{size=\fontsize{13}{16}\selectfont}
\setbeamersize{text margin left=0.83in,text margin right=0.83in}
\setbeamertemplate{navigation symbols}{}
\setbeamertemplate{frametitle}{\vspace{0.26in}\insertframetitle\par\vspace{0.18in}}
\setbeamertemplate{footline}{\hfill\color{muted}\insertframenumber\,/\,\inserttotalframenumber\hspace{0.83in}\vspace{0.23in}}
\setbeamertemplate{itemize item}{\small\raise0.12ex\hbox{\textbullet}}
\setbeameroption{hide notes}
\newcommand{\E}{\mathbb{E}}
\newcommand{\Prb}{\mathbb{P}}
\newcommand{\ind}{\mathbf{1}}
\newcommand{\quiet}[1]{\par{\fontsize{17}{23}\selectfont\color{muted}#1\par}}
\newcommand{\gap}{\par\vspace{0.23in}}
\hypersetup{pdftitle={CAST: Class slides for causal educators},pdfsubject={Editable draft slides on potential outcomes and survival-effect trajectories}}
\begin{document}
\usebeamerfont{normal text}

\begin{frame}[plain]
\vspace{0.1in}
{\fontsize{20}{26}\selectfont\color{teal}CAST\enspace /\enspace Causal educator kit\par}
\vspace{0.45in}
{\fontsize{49}{56}\selectfont\bfseries One person,\par two possible paths\par}
\vspace{0.4in}
{\fontsize{26}{34}\selectfont Causal effects across time and people\par}
\vspace{0.7in}
\quiet{60--90 minute class \enspace\textperiodcentered\enspace Synthetic examples\newline
Editable draft with a study guide, R lab and demo data}
\note{Open with Hector receiving treatment and improving. Ask what would have happened without treatment. The kit follows the existing interactive story. This is a teaching example, not patient evidence. Source and implementation: https://github.com/ishuryak/cast_demo_website.}
\end{frame}

\begin{frame}{The missing comparison}
Hector receives treatment and improves.
\gap
What would have happened to \emph{Hector}, at the \emph{same time},
without treatment?
\par\vspace{0.4in}
{\fontsize{30}{38}\selectfont\color{teal}One assigned path can be observed.\par
The alternative remains unobserved.\par}
\gap
\quiet{Before versus after compares two times on one path.\newline
Treatment versus control compares two potential paths at one time.}
\note{Invite learners to articulate the fundamental problem before showing symbols. In the opening artwork, lower symptom burden represents improvement; the rest of this class uses survival, where higher is better. Censoring can obscure an outcome even on the assigned path. Randomization identifies averages under assumptions; it does not reveal both outcomes for one person. Foundation: Hernan and Robins, Causal Inference: What If, chapters 1--3, https://miguelhernan.org/whatifbook.}
\end{frame}

\begin{frame}{A survival snapshot}
At horizon $t$, let $A^w(t)=\ind\{T(w)>t\}$ mean survival under
baseline treatment choice $w$.
\gap
\[
 \Delta(t)=\Prb\{T(1)>t\}-\Prb\{T(0)>t\}
\]
\gap
{\color{teal}$+0.10$ means a 10 percentage-point survival advantage.}
\gap
\quiet{Specify the population, treatment choices and horizon.\newline
This is a probability difference, not a hazard ratio or a relative increase.}
\note{T(w) denotes a potential event time. A_i^w(t) is 0 or 1, while population survival probabilities range from 0 to 1. Positive values favor treatment on the survival scale. Ask a learner to translate 0.08 into percentage points. grf survival-probability target: https://grf-labs.github.io/grf/reference/causal_survival_forest.html.}
\end{frame}

\begin{frame}{The individual and the average}
\renewcommand{\arraystretch}{1.7}
\begin{tabular}{@{}p{2.4in}p{7.7in}@{}}
\textcolor{teal}{One person} & $\delta_i(t)=A_i^1(t)-A_i^0(t)$ \\
Realized sample & $\displaystyle\Delta_{\mathrm{SATE}}(t)=\frac1n\sum_i\delta_i(t)$ \\
Population & $\displaystyle\Delta(t)=\E[\delta(t)]$ \\
Similar profiles & $\displaystyle\tau(t,x)=\E[\delta(t)\mid X=x]$ \\
\end{tabular}\par
\gap
\quiet{A conditional average for people like Hector is still an average.\newline
It does not identify his missing potential outcome.}
\note{Use guide page 2 to distinguish the realized SATE from an average of conditional survival probabilities over observed profiles. The simulator's oracle is the latter, not paired realized potential outcomes. Discuss exercise A: H-L=2 in ten people could mean two helped and none harmed, or four helped and two harmed. Source: Hernan and Robins, https://miguelhernan.org/whatifbook.}
\end{frame}

\begin{frame}{A fair comparison}
Healthier people may be more likely to receive treatment\newline
and more likely to survive.
\gap
Adjustment needs measured common causes, treatment\newline
overlap and a well-defined treatment comparison.
\gap
{\color{teal}Good prediction does not establish causal identification.}
\gap
\quiet{Exchangeability: $\{T(0),T(1)\}\perp W\mid X$\newline
Overlap: $0<\Prb(W=1\mid X=x)<1$ for relevant profiles.}
\note{Ask learners to name a plausible unmeasured common cause in their own study. Neither a flexible forest nor a narrow interval repairs hidden confounding. Consistency, absent or explicitly modeled interference, and a target-population argument are also needed. Guide page 3 gives the fuller assumptions. Sources: Hernan and Robins; Cui et al. (2023), https://doi.org/10.1093/jrsssb/qkac001.}
\end{frame}

\begin{frame}{Censoring hides part of the assigned path}
We record an event time or the end of observed follow-up:
\gap
\[
 Y=\min(T,C),\qquad D=\ind\{T\le C\}.
\]
\gap
If follow-up ends before $t$, survival at $t$ may be unknown.
\gap
\quiet{Censoring needs assumptions and adequate follow-up support.\newline
It is separate from the missing counterfactual.}
\note{For example, someone censored at month 20 cannot simply be counted as alive or dead at month 36. A sufficient working condition is T independent of C conditional on X,W, together with censoring positivity through the desired horizons. A late horizon can be unsupported even when treatment overlap is good. Source: Cui et al. (2023), https://doi.org/10.1093/jrsssb/qkac001; grf reference.}
\end{frame}

\begin{frame}{The effect can change over follow-up}
Causal survival forests (CSF) estimate survival effects.\newline
The website displays cohort averages at five horizons.
\gap
CAST fits a trajectory to those estimates:
\[
 \widehat\Delta_{\mathrm{CAST}}(t)=\widehat\beta_0+
 \widehat\beta_1t+\widehat\beta_2t^2.
\]
\gap
\quiet{The fitted line adds a model of the trajectory.\newline
Additional drawing points add no observations. Keep the CSF points visible.}
\note{CSF can estimate conditional averages; the website's displayed target is the cohort average. The demo fits a quadratic; the CAST preprint also discusses other trajectory choices. The horizon grid is not a binning of each patient's event time. Open the interactive chart here if available. Sources: Yang et al., https://arxiv.org/abs/2505.06367; exact demo implementation, https://github.com/ishuryak/cast_demo_website/blob/43a00101100801eb6c048800d24788455fac0f8b/R/cast_core.R.}
\end{frame}

\begin{frame}{Read uncertainty at the right level}
The same people contribute at several horizons,\newline
so the estimates are correlated.
\gap
The demo carries that covariance into a\newline
\alert{95\% pointwise band} around the fitted quadratic.
\gap
\quiet{Pointwise is not simultaneous coverage of the whole curve.\newline
The band does not capture all fit-selection or model-misspecification uncertainty.}
\note{Optional math is on guide page 4. The demo estimates covariance from patient-aligned influence scores, applies shrinkage and selects GLS when conditioning and fit checks pass, with WLS fallback. The score construction and selection are implementation choices; do not attribute all of them to the original CAST preprint. Exact source: https://github.com/ishuryak/cast_demo_website/blob/43a00101100801eb6c048800d24788455fac0f8b/R/cast_core.R. A declining positive survival difference still favors treatment at that horizon. A peak does not identify when to stop treatment.}
\end{frame}

\begin{frame}{A smooth fit can miss the trajectory}
\quiet{Frozen reversal scenario \enspace\textperiodcentered\enspace 2,000 synthetic people\newline
Measured confounding 1; hidden confounding 0. Differences in pp.}
\vspace{0.2in}
\centering
\renewcommand{\arraystretch}{1}
\begin{tabular}{@{}>{\rule[-9pt]{0pt}{37pt}}rrrrr@{}}
\toprule
Month & Unadjusted & CSF & CAST & Oracle \\
\midrule
12  & 13.49 & 6.51 & 6.39 & 7.51 \\
36  & 32.35 & 18.09 & 10.34 & 16.38 \\
60  & 28.56 & 11.42 & 9.96 & 9.15 \\
84  & 14.51 & $-0.12$ & 5.28 & $-1.59$ \\
108 & 6.81 & $-4.38$ & $-3.72$ & $-5.60$ \\
\bottomrule
\end{tabular}\par
\vspace{0.25in}
\raggedright
\quiet{At which horizon is the CAST--oracle gap largest?}
\note{Answer: 84 months, 6.87 pp, versus 6.04 pp at 36 months. At 108 months the naive comparison remains positive while CSF, CAST and the oracle are negative. An oracle here is the known generating model's survival-probability contrast averaged over the cohort profiles, not each person's realized causal outcome. Source: the frozen scenarios.json export, scenario reversal_conf1.00_unmeas0.00, reproduced in data/frozen-horizons.csv in this kit. Values are shown to two decimals in percentage points. This is not the kit's separate 600-person cohort.}
\end{frame}

\begin{frame}{For whom, and for which question?}
An effect may vary with baseline profile $x$:
\[
 \tau(t,x)=\E[A^1(t)-A^0(t)\mid X=x].
\]
\gap
Name the outcome scale, horizon and population\newline
before describing heterogeneity.
\gap
\quiet{These graphics show cohort averages, not subgroup estimates.\newline
The effect of baseline treatment does not answer when to switch or stop it.}
\note{A common relative hazard effect can coexist with different absolute survival benefits across risk profiles. This class uses absolute survival-probability differences. Following a baseline contrast over time does not estimate a dynamic treatment regime. Choosing whether and when to change treatment requires a different estimand and identification argument. Source: grf CSF reference, https://grf-labs.github.io/grf/reference/causal_survival_forest.html.}
\end{frame}

\begin{frame}{The class lab}
Inspect the supplied \alert{600-person synthetic cohort}.\newline
Run \texttt{class-lab.R}, or use the frozen results.
\gap
Compare five horizon estimates with the separate answer key.\newline
Explain one discrepancy and one assumption the data cannot prove.
\gap
\quiet{The small lab is a different sample from the website's 2,000-person export.\newline
Record the target, horizon, units and uncertainty in your interpretation.}
\note{The cohort uses seed 20260905, reversal effects, measured confounding 1 and no hidden confounding. It records age, stage, ps, comorb, smoke, sex, ethnicity, W, Y and D. The answer key averages known generating survival probabilities for those same profiles. Do not supply answer-key values as fitting inputs. Code runs Igor's implementation; consult the kit README for packages and exact invocation. Compare sample estimates to the answer key without expecting equality. For a no-compute class, use data/frozen-horizons.csv and exercise B.}
\end{frame}

\begin{frame}{Return to the causal roadmap}
Define the question and target population.\newline
Connect the data through explicit assumptions.\newline
Estimate, then interpret with uncertainty.
\gap
{\color{teal}CSF and CAST belong inside that process.}
\gap
\quiet{Igor Shuryak's \href{https://github.com/ishuryak/cast_demo_website}{CAST repository}: R implementation and methods\newline
\href{https://miguelhernan.org/whatifbook}{Hern\'an \& Robins}: causal foundations \enspace\textperiodcentered\enspace
\href{https://doi.org/10.1093/jrsssb/qkac001}{Cui et al.}: CSF\newline
\href{https://arxiv.org/abs/2505.06367}{Yang et al.}: CAST \enspace\textperiodcentered\enspace
\href{https://doi.org/10.1097/EDE.0000000000000078}{Petersen \& van der Laan}: causal roadmap\newline
\href{https://navigator.tao-rwd.com/}{Tao of RWD Navigator}: a place to develop the study question}
\note{Close by asking learners to write a target question for their own work and name a reason this method might not fit it. This roadmap is a teaching adaptation, not the source paper's exact step list. The study guide contains vocabulary, equations, exercises and answer notes. Slides and guide are editable LaTeX. Sources: Petersen and van der Laan (2014), doi:10.1097/EDE.0000000000000078; Cui et al. (2023), doi:10.1093/jrsssb/qkac001; Yang et al. (2025), arXiv:2505.06367; Hernan and Robins (2020), Causal Inference: What If.}
\end{frame}
\end{document}
